\chapter{Galois 理论初步}

上一章已经介绍了分裂域、可分扩张与正规扩张。本章要说明：对于一个有限可分正规扩张 $L/K$，群
\[
\Gal(L/K)=\Aut_K(L)
\]
能够完整编码这个扩张的内部结构。Galois 理论的基本定理把中间域与 Galois 群的子群一一对应；而群的可解性又控制多项式能否由根式求解。最后，我们简要概览绝对 Galois 群、$\ell$-进表示与模形式之间的联系。

本章默认所有扩张都嵌入某个固定的代数闭包中。若无特别说明，扩张均指有限扩张。

\section{Galois 群的定义}

\subsection*{定义}

\begin{definition}
设 $L/K$ 是域扩张。称
\[
\Gal(L/K)=\Aut_K(L)=\{\sigma:L\to L \mid \sigma \text{ 是域自同构且 } \sigma|_K=\Id_K\}
\]
为扩张 $L/K$ 的 \emph{Galois 群}。
\end{definition}

\begin{proposition}
在映射复合下，$\Gal(L/K)$ 构成一个群。
\end{proposition}

\begin{proof}
恒等映射显然在其中。若 $\sigma,\tau$ 都固定 $K$，则 $\sigma\circ\tau$ 仍是 $L$ 的域自同构，且对任意 $a\in K$ 有
\[
(\sigma\circ\tau)(a)=\sigma(a)=a。
\]
故 $\sigma\circ\tau$ 也固定 $K$。又若 $\sigma\in\Gal(L/K)$，则 $\sigma^{-1}$ 仍是域自同构，并满足 $\sigma^{-1}(a)=a$ 对所有 $a\in K$ 成立，因此逆元也在其中。故它构成一个群。
\end{proof}

\begin{remark}
Galois 群本质上描述的是扩张中的“对称性”：它保留底域中的系数不变，只在上层域中重新排列共轭元素，尤其是重新排列多项式的根。
\end{remark}

\subsection*{基本例子}

\begin{example}[二次扩张]
若 $d\in\Q$ 不是有理数平方，则
\[
\Gal(\Q(\sqrt d)/\Q)\cong \Z/2\Z。
\]
特别地，$\Gal(\Q(\sqrt2)/\Q)\cong \Z/2\Z$。
\end{example}

\begin{proof}
扩张 $\Q(\sqrt d)/\Q$ 的次数为 $2$，所以任何 $\Q$-自同构都由 $\sqrt d$ 的像唯一决定。又
\[
(\sigma(\sqrt d))^2=\sigma(d)=d，
\]
故 $\sigma(\sqrt d)$ 只能是 $\sqrt d$ 或 $-\sqrt d$。因此至多有两个自同构。另一方面，映射
\[
\tau(a+b\sqrt d)=a-b\sqrt d,
\qquad (a,b\in\Q)
\]
确实给出一个非平凡的 $\Q$-自同构，所以 Galois 群恰有两个元素，与 $\Z/2\Z$ 同构。
\end{proof}

\begin{example}[双二次扩张]
设
\[
L=\Q(\sqrt2,\sqrt3)。
\]
则
\[
\Gal(L/\Q)\cong (\Z/2\Z)^2。
\]
\end{example}

\begin{proof}
任一 $\Q$-自同构必须把 $\sqrt2$ 送到 $\pm\sqrt2$，把 $\sqrt3$ 送到 $\pm\sqrt3$，故至多有四个自同构。而四种符号组合都确实给出自同构，因此群恰有四个元素。它由
\[
\sigma(\sqrt2)=-\sqrt2,
\quad
\sigma(\sqrt3)=\sqrt3；
\]
\[
\tau(\sqrt2)=\sqrt2,
\quad
\tau(\sqrt3)=-\sqrt3
\]
生成，且 $\sigma^2=\tau^2=1$、$\sigma\tau=\tau\sigma$，因此 Galois 群同构于 Klein 四元群 $(\Z/2\Z)^2$。
\end{proof}

\begin{example}[有限域]
对任意素数 $p$ 与正整数 $n$，有
\[
\Gal(\F_{p^n}/\F_p)\cong \Z/n\Z，
\]
其生成元可取 Frobenius 自同构
\[
\Frob_p:x\mapsto x^p。
\]
\end{example}

\begin{proof}
在特征 $p$ 下，$(x+y)^p=x^p+y^p$，因此映射 $x\mapsto x^p$ 是域同态；有限域上的单射自动满射，所以它是 $\F_{p^n}$ 的自同构。又对所有 $a\in\F_p$，有 $a^p=a$，故此自同构固定 $\F_p$。

对任意 $x\in\F_{p^n}$，有 $x^{p^n}=x$，故 $\Frob_p^n=\Id$。若存在 $0<m<n$ 使 $\Frob_p^m=\Id$，则 $\F_{p^n}$ 的所有元素都满足 $x^{p^m}=x$，于是多项式 $X^{p^m}-X$ 有至少 $p^n$ 个根，这与其次数 $p^m<p^n$ 矛盾。因此 $\Frob_p$ 的阶为 $n$。

再取 $\alpha\in\F_{p^n}$ 使 $\F_{p^n}=\F_p(\alpha)$。任何 $\F_p$-自同构由 $\alpha$ 的像唯一决定，而 $\alpha$ 的共轭根正是
\[
\alpha,\alpha^p,\alpha^{p^2},\dots,\alpha^{p^{n-1}}。
\]
故每个 $\F_p$-自同构都是某个 $\Frob_p^i$。于是 Galois 群由 $\Frob_p$ 生成，因而是阶为 $n$ 的循环群。
\end{proof}

\begin{example}[分圆域]
若 $\zeta_n$ 是原始 $n$ 次单位根，则对每个与 $n$ 互素的整数 $a$，赋值
\[
\sigma_a(\zeta_n)=\zeta_n^a
\]
给出一个 $\Q$-自同构。因此存在单射
\[
(\Z/n\Z)^\times \hookrightarrow \Gal(\Q(\zeta_n)/\Q)。
\]
在 12.4 中我们将证明这其实是同构。
\end{example}

\begin{proof}
因为 $a$ 与 $n$ 互素时，$\zeta_n^a$ 仍是原始 $n$ 次单位根，所以它与 $\zeta_n$ 满足同一个最小多项式，于是把 $\zeta_n$ 送到 $\zeta_n^a$ 可唯一延拓为一个 $\Q$-嵌入。有限扩张上的嵌入必为自同构。若 $\sigma_a=\sigma_b$，则 $\zeta_n^a=\zeta_n^b$，从而 $n\mid(a-b)$，故得到单射。
\end{proof}

\begin{example}[三次分裂域]
多项式 $x^3-2$ 的分裂域为
\[
L=\Q(\sqrt[3]{2},\omega)，
\qquad
\omega=\frac{-1+\sqrt{-3}}{2}。
\]
并且
\[
\Gal(L/\Q)\cong S_3。
\]
\end{example}

\begin{proof}
由 Eisenstein 判别法，$x^3-2$ 在 $\Q$ 上不可约，因此 $[\Q(\sqrt[3]{2}):\Q]=3$。其三个根为
\[
\sqrt[3]{2},
\quad
\omega\sqrt[3]{2},
\quad
\omega^2\sqrt[3]{2}。
\]
所以分裂域必须同时加入 $\omega$。又 $\Q(\sqrt[3]{2})\subseteq\R$ 而 $\Q(\omega)\not\subseteq\R$，故交为 $\Q$，于是 $[L:\Q]=6$。

任一 $\Q$-自同构必须把 $\sqrt[3]{2}$ 送到上述三个根之一，把 $\omega$ 送到 $\omega$ 或 $\omega^2$，所以至多有六个自同构。另一方面，映射
\[
\sigma(\sqrt[3]{2})=\omega\sqrt[3]{2},
\qquad
\sigma(\omega)=\omega；
\]
\[
\tau(\sqrt[3]{2})=\sqrt[3]{2},
\qquad
\tau(\omega)=\omega^2
\]
都给出自同构，且满足 $\sigma^3=\tau^2=1$ 与 $\tau\sigma\tau=\sigma^{-1}$，故它们生成一个与 $S_3$ 同构的六阶群。因此整个 Galois 群就是 $S_3$。
\end{proof}

\section{Galois 对应}

本节证明有限 Galois 扩张的基本定理。思想非常朴素：给定中间域 $E$，考虑固定 $E$ 的自同构；给定子群 $H$，考虑被 $H$ 全部元素固定的元素。结论是这两个过程互为逆运算。

\subsection*{固定域}

\begin{definition}
设 $L$ 是域，$G\le\Aut(L)$。定义固定域
\[
L^G=\{x\in L \mid \sigma(x)=x\text{ 对一切 }\sigma\in G\}。
\]
若 $K\subseteq E\subseteq L$，则称 $E$ 为扩张 $L/K$ 的中间域。
\end{definition}

\begin{lemma}
设 $L/K$ 是有限 Galois 扩张，$E$ 为中间域。则 $L/E$ 仍是有限 Galois 扩张，并且
\[
|\Gal(L/E)|=[L:E]。
\]
\end{lemma}

\begin{proof}
因为 $L/K$ 可分，所以中间扩张 $L/E$ 也可分。若 $f(x)\in E[x]$ 在 $L$ 中有一个根，那么该根在 $K$ 上的最小多项式在 $L$ 中完全分裂；于是 $f$ 的全部共轭根仍在 $L$ 中，所以 $L/E$ 正规。故 $L/E$ 为 Galois 扩张。

有限可分扩张的 $E$-嵌入个数等于扩张次数 $[L:E]$；正规性保证这些嵌入的像都落在 $L$ 内，因此恰好给出全部 $E$-自同构，故结论成立。
\end{proof}

\begin{theorem}[固定子群次数公式]
设 $L/K$ 为有限 Galois 扩张，$H\le G=\Gal(L/K)$。记 $E=L^H$。则
\[
[L:E]=|H|，
\qquad
\Gal(L/E)=H。
\]
\end{theorem}

\begin{proof}
首先，$H$ 中每个元素都固定 $E=L^H$，所以 $H\subseteq\Gal(L/E)$。由上引理，
\[
|H|\le |\Gal(L/E)|=[L:E]。
\]

接着证明反向不等式。因为 $L/E$ 是有限可分扩张，由原始元定理可取 $\alpha\in L$ 使 $L=E(\alpha)$。考虑轨道多项式
\[
P(T)=\prod_{\sigma\in H}(T-\sigma(\alpha))。
\]
对任意 $\tau\in H$，集合 $\{\tau\sigma(\alpha):\sigma\in H\}$ 与 $\{\sigma(\alpha):\sigma\in H\}$ 相同，因此 $\tau$ 固定 $P(T)$ 的系数。故 $P(T)\in E[T]$。于是 $\alpha$ 在 $E$ 上的最小多项式整除 $P(T)$，从而
\[
[L:E]=\deg m_{\alpha,E}(T)\le\deg P(T)=|H|。
\]
结合前面的不等式，得 $[L:E]=|H|$。于是
\[
|\Gal(L/E)|=[L:E]=|H|。
\]
而 $H\subseteq\Gal(L/E)$，所以只能有 $\Gal(L/E)=H$。
\end{proof}

\begin{theorem}[有限 Galois 扩张的基本定理]
设 $L/K$ 为有限 Galois 扩张，记 $G=\Gal(L/K)$。则映射
\[
E\mapsto\Gal(L/E)
\]
在中间域集合
\[
\{E\mid K\subseteq E\subseteq L\}
\]
与子群集合
\[
\{H\mid H\le G\}
\]
之间给出一个反向包含双射，其逆映射为
\[
H\mapsto L^H。
\]
此外，对任意中间域 $E$ 都有
\[
[L:E]=|\Gal(L/E)|，
\qquad
[E:K]=[G:\Gal(L/E)]。
\]
\end{theorem}

\begin{proof}
先取中间域 $E$。令 $H=\Gal(L/E)$。由固定子群次数公式，
\[
[L:L^H]=|H|=|\Gal(L/E)|=[L:E]。
\]
又显然 $E\subseteq L^H$，故只能有 $E=L^H$。

再取任意子群 $H\le G$。令 $E=L^H$。由固定子群次数公式，立即得
\[
\Gal(L/E)=H。
\]
因此这两个映射互为逆映射，从而构成双射。

若 $E_1\subseteq E_2$，则固定较大域的自同构更少，因此
\[
\Gal(L/E_2)\subseteq\Gal(L/E_1)。
\]
这说明包含关系反向。最后，由塔式公式和 $|G|=[L:K]$ 得
\[
[E:K]=\frac{[L:K]}{[L:E]}=\frac{|G|}{|\Gal(L/E)|}=[G:\Gal(L/E)]。
\]
\end{proof}

\begin{proposition}[正规子群与正规中间扩张]
设 $L/K$ 为有限 Galois 扩张，$E$ 为中间域，$H=\Gal(L/E)$。则下列条件等价：
\begin{enumerate}[label=\normalfont(\arabic*), leftmargin=2.5em]
\item $E/K$ 是 Galois 扩张；
\item $H\lhd\Gal(L/K)$。
\end{enumerate}
在这时，限制映射
\[
\Gal(L/K)\to\Gal(E/K)，
\qquad
\sigma\mapsto\sigma|_E
\]
是满同态，且有
\[
\Gal(E/K)\cong \Gal(L/K)/H。
\]
\end{proposition}

\begin{proof}
先证 $(2)\Rightarrow(1)$。取 $\tau\in\Gal(L/K)$ 与 $x\in E=L^H$。若 $\sigma\in H$，则由正规性知 $\tau^{-1}\sigma\tau\in H$，于是
\[
\sigma(\tau(x))=\tau((\tau^{-1}\sigma\tau)(x))=\tau(x)。
\]
故 $\tau(x)$ 被 $H$ 中全部元素固定，故 $\tau(x)\in E$。因此每个 $K$-自同构都稳定 $E$，从而 $E/K$ 正规；又中间扩张仍可分，所以 $E/K$ 为 Galois 扩张。

再证 $(1)\Rightarrow(2)$。任取 $\tau\in\Gal(L/K)$ 与 $\sigma\in H$。若 $x\in E$，由于 $E/K$ 正规，$\tau(x)\in E$；又 $\sigma$ 固定 $E$，所以
\[
(\tau^{-1}\sigma\tau)(x)=\tau^{-1}(\sigma(\tau(x)))=x。
\]
故 $\tau^{-1}\sigma\tau\in H$，即 $H$ 为正规子群。

若两条件成立，则限制映射的核正是固定 $E$ 的自同构群，即 $H$。为证满射，任取 $\rho\in\Gal(E/K)$。因为 $L/E$ 为有限可分扩张，$\rho$ 可以延拓为一个 $K$-嵌入 $L\hookrightarrow\bar K$；又由于 $L/K$ 正规，其像仍在 $L$ 中，因此该延拓实际上属于 $\Gal(L/K)$。于是限制映射满，由第一同构定理即得结论。
\end{proof}

\subsection*{格图示例}

\begin{example}[双二次扩张]
对扩张
\[
L=\Q(\sqrt2,\sqrt3)，
\qquad
G=\Gal(L/\Q)\cong(\Z/2\Z)^2，
\]
中间域格与子群格如下。
\end{example}

\begin{center}
\begin{tikzpicture}[node distance=1.6cm and 2.2cm]
\node (L) {$L=\Q(\sqrt2,\sqrt3)$};
\node (E1) [below left=of L] {$\Q(\sqrt2)$};
\node (E2) [below=of L] {$\Q(\sqrt6)$};
\node (E3) [below right=of L] {$\Q(\sqrt3)$};
\node (K) [below=of E2] {$\Q$};
\draw (L)--(E1) (L)--(E2) (L)--(E3);
\draw (E1)--(K) (E2)--(K) (E3)--(K);
\end{tikzpicture}
\hspace{2cm}
\begin{tikzpicture}[node distance=1.6cm and 2.2cm]
\node (T) {$\{1\}$};
\node (H1) [below left=of T] {$\langle \sigma \rangle$};
\node (H2) [below=of T] {$\langle \sigma\tau \rangle$};
\node (H3) [below right=of T] {$\langle \tau \rangle$};
\node (B) [below=of H2] {$G$};
\draw (T)--(H1) (T)--(H2) (T)--(H3);
\draw (H1)--(B) (H2)--(B) (H3)--(B);
\end{tikzpicture}
\end{center}

左图中的三个二次子域分别对应右图中的三个二阶子群，而且包含关系完全反向。

\begin{example}[三次分裂域]
若 $L=\Q(\sqrt[3]{2},\omega)$，则 Galois 群同构于 $S_3$。因为 $S_3$ 具有一个三阶正规子群和三个二阶子群，所以 $L/\Q$ 恰有一个二次中间域 $\Q(\omega)$，以及三个互相共轭的三次中间域。只有二次中间域对应正规子群，因此只有它是 $\Q$ 上的 Galois 扩张。
\end{example}

\section{可解群与根式可解性}

\subsection*{可解群}

\begin{definition}
设 $G$ 为群。其导出列定义为
\[
G^{(0)}=G，
\qquad
G^{(i+1)}=[G^{(i)},G^{(i)}]。
\]
若存在 $m\ge0$ 使得 $G^{(m)}=1$，则称 $G$ 为可解群。
\end{definition}

\begin{proposition}
下列命题成立：
\begin{enumerate}[label=\normalfont(\arabic*), leftmargin=2.5em]
\item 阿贝尔群都是可解群；
\item 可解群的子群与商群仍可解；
\item 若 $N\lhd G$ 且 $N$ 与 $G/N$ 都可解，则 $G$ 可解。
\end{enumerate}
\end{proposition}

\begin{proof}
阿贝尔群的交换子群平凡，因此可解。若 $H\le G$，则 $H^{(i)}\subseteq G^{(i)}$ 对所有 $i$ 成立，所以 $G$ 可解推出 $H$ 可解。对商群，设 $\pi:G\twoheadrightarrow G/N$ 为自然投影，则 $\pi(G^{(i)})=(G/N)^{(i)}$，所以 $G$ 可解推出 $G/N$ 可解。

若 $N^{(r)}=1$ 且 $(G/N)^{(s)}=1$，则由投影映射可知 $G^{(s)}\subseteq N$，因而
\[
G^{(s+r)}=(G^{(s)})^{(r)}\subseteq N^{(r)}=1。
\]
故 $G$ 可解。
\end{proof}

\begin{example}
二面体群 $D_{2n}$ 可解，因为它有正规循环子群 $C_n$，而商群同构于 $C_2$。另一方面，
\[
1\lhd V_4\lhd A_4\lhd S_4
\]
给出 $S_4$ 的一个商列，其各商群依次为 $C_2,C_3,C_2$，所以 $S_4$ 也是可解群。
\end{example}

\subsection*{$S_n$ 在 $n\ge5$ 时不可解}

\begin{proposition}
对一切 $n\ge2$，有
\[
[S_n,S_n]=A_n。
\]
\end{proposition}

\begin{proof}
符号同态 $\operatorname{sgn}:S_n\to\{\pm1\}$ 的核是 $A_n$，所以交换子群包含于 $A_n$。另一方面，$A_n$ 由 3-循环生成，而
\[
(1\,2\,3)=(1\,2)(2\,3)(1\,2)(2\,3)^{-1}
\]
是一个交换子，因此所有 3-循环都属于 $[S_n,S_n]$，从而 $A_n\subseteq[S_n,S_n]$。故等号成立。
\end{proof}

\begin{theorem}
当 $n\ge5$ 时，$S_n$ 不是可解群。
\end{theorem}

\begin{proof}
已知 $A_n$ 在 $n\ge5$ 时是非阿贝尔单群。由上命题，$S_n'=A_n$。又因为 $A_n$ 非阿贝尔且单，故其交换子群 $A_n'$ 是一个非平凡正规子群，只能等于 $A_n$ 本身。于是从第二步开始，$S_n$ 的导出列恒等于 $A_n$，不可能下降到平凡群。因此 $S_n$ 不可解。
\end{proof}

\begin{remark}
这解释了为什么三次与四次一般方程存在统一的根式公式，而五次一般方程不存在：本质区别就在于 $S_3,S_4$ 可解，而 $S_n$ 在 $n\ge5$ 时不可解。
\end{remark}

\subsection*{根式可解性}

\begin{definition}
设 $K$ 的特征为 $0$。若存在域塔
\[
K=K_0\subseteq K_1\subseteq\cdots\subseteq K_r
\]
使得对每个 $i$，都有某个元素 $\alpha_i\in K_i$ 与正整数 $m_i$ 满足
\[
K_i=K_{i-1}(\alpha_i)，
\qquad
\alpha_i^{m_i}\in K_{i-1}，
\]
则称 $K_r/K$ 为根式扩张。若多项式 $f\in K[x]$ 的全部根包含于某个根式扩张中，则称 $f$ 在 $K$ 上可由根式求解。
\end{definition}

\begin{proposition}
若 $K$ 含有原始 $m$ 次单位根 $\zeta_m$，并且 $L$ 是多项式 $x^m-a$ 在 $K$ 上的分裂域，则 $\Gal(L/K)$ 为阿贝尔群。
\end{proposition}

\begin{proof}
取根 $\alpha$ 使 $\alpha^m=a$。由于 $K$ 已含有全部 $m$ 次单位根，$x^m-a$ 的全部根为
\[
\alpha,\zeta_m\alpha,\dots,\zeta_m^{m-1}\alpha，
\]
故分裂域就是 $L=K(\alpha)$。任一 $\sigma\in\Gal(L/K)$ 满足
\[
(\sigma(\alpha))^m=\sigma(a)=a，
\]
因此 $\sigma(\alpha)=\zeta_m^j\alpha$ 对某个 $j$ 成立。于是映射
\[
\Gal(L/K)\to\mu_m，
\qquad
\sigma\mapsto\frac{\sigma(\alpha)}{\alpha}
\]
是单射，故 $\Gal(L/K)$ 同构于循环群 $\mu_m$ 的某个子群，特别是阿贝尔群。
\end{proof}

\begin{theorem}
若特征为 $0$ 的域 $K$ 上多项式 $f$ 可由根式求解，则其分裂域的 Galois 群是可解群。
\end{theorem}

\begin{proof}
设 $f$ 的全部根落在某个根式扩张塔的顶层域中。把每一步所需的单位根也一并加入塔中；这些分圆扩张的 Galois 群是阿贝尔群，因此是可解的。于是可以假设每一步 $K_i/K_{i-1}$ 都由加入一个满足 $\alpha_i^{m_i}\in K_{i-1}$ 的根式元素得到，并且底域已经含有原始 $m_i$ 次单位根。

对每一步，上一命题说明相应的 Galois 闭包相对于前一步的商群是阿贝尔的。于是逐层应用“可解正规子群与可解商群推出整体可解”的判据，可得最上层 Galois 闭包对 $K$ 的 Galois 群可解。分裂域对应于这个 Galois 群的一个商群，因此也可解。
\end{proof}

\begin{theorem}[Abel--Ruffini 定理]
当次数 $n\ge5$ 时，一般 $n$ 次多项式不能由根式统一求解。
\end{theorem}

\begin{proof}[证明思路]
把一般 $n$ 次多项式看作
\[
X^n+t_1X^{n-1}+\cdots+t_n
\]
在有理函数域 $\Q(t_1,\dots,t_n)$ 上的多项式。它的 Galois 群是 $S_n$。若 $n\ge5$，则前定理表明 $S_n$ 不可解，所以一般方程不可能由根式统一求解。
\end{proof}

\begin{remark}
Abel--Ruffini 定理说的是“一般方程”不可解，而不是每个五次方程都不可解。某些特殊五次方程的 Galois 群仍可能是可解群，从而依然可以由根式求解。
\end{remark}

\section{分圆域与分圆多项式}

分圆域是最重要的一类显式 Galois 扩张。它们一方面提供了大量可计算的例子，另一方面又是阿贝尔扩张理论的模型。

\subsection*{分圆多项式的定义}

\begin{definition}
设 $n\ge1$。定义第 $n$ 个分圆多项式为
\[
\Phi_n(x)=\prod_{1\le a\le n,\,(a,n)=1}(x-\zeta_n^a)，
\]
其中 $\zeta_n$ 为任意一个原始 $n$ 次单位根。
\end{definition}

\begin{proposition}
对任意 $n\ge1$，有
\[
x^n-1=\prod_{d\mid n}\Phi_d(x)。
\]
\end{proposition}

\begin{proof}
$ x^n-1 $ 的根恰好是所有 $n$ 次单位根。每个 $n$ 次单位根的阶是某个除数 $d\mid n$，并且它恰是某个原始 $d$ 次单位根。不同的 $d$ 给出互不相交的根集合，而这些集合的并恰好是 $x^n-1$ 的全部根。由于各因子都是首一多项式，所以它们的乘积就是 $x^n-1$。
\end{proof}

\begin{corollary}
有 Möbius 反演公式
\[
\Phi_n(x)=\prod_{d\mid n}(x^d-1)^{\mu(n/d)}。
\]
\end{corollary}

\begin{proof}
由上一命题在除数格上施行 Möbius 反演即可。
\end{proof}

\begin{example}
前几个分圆多项式为
\[
\Phi_1(x)=x-1,
\quad
\Phi_2(x)=x+1,
\quad
\Phi_3(x)=x^2+x+1,
\quad
\Phi_4(x)=x^2+1；
\]
\[
\Phi_5(x)=x^4+x^3+x^2+x+1,
\quad
\Phi_6(x)=x^2-x+1,
\quad
\Phi_8(x)=x^4+1。
\]
\end{example}

\subsection*{不可约性}

\begin{theorem}
对任意正整数 $n$，分圆多项式 $\Phi_n(x)$ 在 $\Q$ 上不可约。
\end{theorem}

\begin{proof}
设 $f(x)\in\Z[x]$ 是 $\Phi_n(x)$ 的一个首一不可约因子，取其根 $\alpha$。因为 $\alpha$ 是原始 $n$ 次单位根，我们只需证明：若素数 $p$ 不整除 $n$，则 $\alpha^p$ 也是 $f$ 的根。如此便可推出与 $n$ 互素的所有指数幂 $\alpha^a$ 都是 $f$ 的根，从而 $f$ 包含全部原始 $n$ 次单位根，只能等于 $\Phi_n(x)$。

反设存在素数 $p$ 满足 $p\nmid n$ 且 $\alpha^p$ 不是 $f$ 的根。取 $\Phi_n(x)$ 的另一个首一不可约因子 $g(x)$ 使得 $g(\alpha^p)=0$。于是 $\alpha$ 是 $g(x^p)$ 的根，所以 $f(x)\mid g(x^p)$。由 Gauss 引理，此整除关系可视为在 $\Z[x]$ 中成立。

模 $p$ 约化后，记约化用横线表示，则
\[
\bar f(x)\mid\overline{g(x^p)}=\bar g(x)^p。
\]
由于 $\bar f(x)$ 在 $\F_p[x]$ 中不可约，故必有 $\bar f(x)\mid\bar g(x)$。这意味着 $\bar f$ 与 $\bar g$ 在某个代数闭包中有公共根，于是 $\overline{\Phi_n(x)}$ 有重根。

另一方面，当 $p\nmid n$ 时，$x^n-1$ 在 $\F_p[x]$ 中的导数为 $nx^{n-1}$，与 $x^n-1$ 互素，所以 $x^n-1$ 没有重根。既然 $\overline{\Phi_n(x)}$ 整除 $x^n-1$，它也不可能有重根，矛盾。

故 $\alpha^p$ 必为 $f$ 的根。由素因子分解可知，对每个与 $n$ 互素的整数 $a$，$\alpha^a$ 都是 $f$ 的根。于是 $f$ 的根正是全部原始 $n$ 次单位根，所以 $f=\Phi_n$，定理得证。
\end{proof}

\begin{corollary}
\[
[\Q(\zeta_n):\Q]=\varphi(n)。
\]
\end{corollary}

\begin{proof}
扩张 $\Q(\zeta_n)/\Q$ 的次数等于 $\zeta_n$ 在 $\Q$ 上最小多项式的次数，而该最小多项式正是 $\Phi_n(x)$，其次数为 $\varphi(n)$。
\end{proof}

\begin{theorem}
对任意正整数 $n$，有自然同构
\[
\Gal(\Q(\zeta_n)/\Q)\cong(\Z/n\Z)^\times。
\]
\end{theorem}

\begin{proof}
在 12.1 中已构造出单射
\[
(\Z/n\Z)^\times\hookrightarrow\Gal(\Q(\zeta_n)/\Q)。
\]
左侧群的阶为 $\varphi(n)$，右侧群作为 Galois 群的阶为
\[
|\Gal(\Q(\zeta_n)/\Q)|=[\Q(\zeta_n):\Q]=\varphi(n)。
\]
故此单射必为双射，从而得到同构。
\end{proof}

\begin{example}
$\Gal(\Q(\zeta_5)/\Q)$ 是四阶循环群；$\Gal(\Q(\zeta_8)/\Q)$ 与 $\Gal(\Q(\zeta_{12})/\Q)$ 都同构于 Klein 四元群。
\end{example}

\subsection*{应用}

\begin{proposition}
若 $p$ 为奇素数，则 $\Gal(\Q(\zeta_p)/\Q)$ 是阶为 $p-1$ 的循环群。因此对每个除数 $d\mid(p-1)$，存在唯一的中间域 $E$ 使得 $[E:\Q]=d$。
\end{proposition}

\begin{proof}
因为 $(\Z/p\Z)^\times\cong\F_p^\times$，而有限域乘法群是循环群，所以分圆域的 Galois 群为循环群。循环群对每个除数恰有一个相应阶的子群，故由 Galois 对应得到唯一的中间域。
\end{proof}

\begin{proposition}[正多边形可作的必要条件]
若正 $n$ 边形可由直尺与圆规作出，则 $\varphi(n)$ 必为 $2$ 的幂。
\end{proposition}

\begin{proof}
正 $n$ 边形可作，当且仅当实数 $\cos(2\pi/n)$ 可由反复二次扩张从 $\Q$ 得到。又
\[
\Q(\cos(2\pi/n))=\Q(\zeta_n+\zeta_n^{-1})
\]
是分圆域的最大实子域，其次数为 $\varphi(n)/2$（当 $n>2$ 时）。若它可由反复二次扩张得到，则该次数只能是 $2$ 的幂，所以 $\varphi(n)$ 也必须是 $2$ 的幂。
\end{proof}

\begin{remark}
更精确的 Gauss--Wantzel 定理说明：正 $n$ 边形可作，当且仅当 $n$ 是 $2$ 的幂与若干互异 Fermat 素数的乘积。这一结论本质上来自分圆域 Galois 群的结构。
\end{remark}

\begin{remark}
Kronecker--Weber 定理断言：$\Q$ 的每个有限阿贝尔扩张都包含在某个分圆域中。它说明分圆域并非若干孤立的例子，而是阿贝尔扩张理论的总模型。
\end{remark}

\section{Galois 表示与模形式（概览）}

这一节不再追求完整证明，而是给出通向现代数论的路线图。前面研究的是有限 Galois 群；更深刻的对象是绝对 Galois 群及其线性表示。

\subsection*{绝对 Galois 群}

\begin{definition}
设 $K$ 是域，$\bar K$ 是其代数闭包。群
\[
G_K=\Gal(\bar K/K)
\]
称为 $K$ 的绝对 Galois 群。特别地，
\[
G_{\Q}=\Gal(\bar{\Q}/\Q)
\]
称为有理数域的绝对 Galois 群。
\end{definition}

\begin{remark}
绝对 Galois 群通常是无限的，并且是一个 profinite 群：
\[
G_K\cong\varprojlim_{L/K\text{ 有限 Galois}}\Gal(L/K)。
\]
因此，对绝对 Galois 群的研究就是对所有有限 Galois 扩张的整体研究。
\end{remark}

\subsection*{$\ell$-进表示}

\begin{definition}
设 $\ell$ 为素数。一个二维 $\ell$-进 Galois 表示是一个连续同态
\[
\rho:G_{\Q}\to\GL_2(\Q_{\ell})
\]
或更一般地
\[
\rho:G_{\Q}\to\GL_2(E)，
\]
其中 $E/\Q_{\ell}$ 为有限扩张。
\end{definition}

\begin{example}
若 $E/\Q$ 是椭圆曲线，则其 $\ell$-进 Tate 模
\[
T_{\ell}(E)=\varprojlim_n E[\ell^n]
\]
是秩为 $2$ 的 $\Z_{\ell}$-模。绝对 Galois 群作用在各层挠点上，取逆极限得到表示
\[
\rho_{E,\ell}:G_{\Q}\to\Aut_{\Z_{\ell}}(T_{\ell}(E))\cong\GL_2(\Z_{\ell})。
\]
\end{example}

\begin{definition}
设 $\rho:G_{\Q}\to\GL_2(E)$，并设 $p\neq\ell$ 为素数。若 $p$ 处的惯性群在该表示下作用平凡，则称 $\rho$ 在 $p$ 处不分歧。此时可以定义 Frobenius 共轭类 $\Frob_p$，并讨论其迹与行列式
\[
\mathrm{tr}(\rho(\Frob_p))，
\qquad
\det(\rho(\Frob_p))。
\]
\end{definition}

\begin{theorem}[Deligne，陈述]
设
\[
f(q)=\sum_{n\ge1}a_nq^n
\]
是权 $k\ge1$、水平 $N$、特征 $\varepsilon$ 的归一化 Hecke 特征模形式。则对每个素数 $\ell$，存在连续半单表示
\[
\rho_{f,\ell}:G_{\Q}\to\GL_2(\bar{\Q}_{\ell})
\]
使得对每个 $p\nmid N\ell$，有
\[
\mathrm{tr}(\rho_{f,\ell}(\Frob_p))=a_p，
\qquad
\det(\rho_{f,\ell}(\Frob_p))=\varepsilon(p)p^{k-1}。
\]
\end{theorem}

\begin{remark}
这条定理把模形式的 Fourier 系数解释为绝对 Galois 群在 Frobenius 元上的迹。它是模形式与算术之间最核心的桥梁之一。
\end{remark}

\begin{theorem}[Serre 模性猜想，现为定理]
每个连续、奇的、不可约的二维模 $\ell$ 表示
\[
\bar\rho:G_{\Q}\to\GL_2(\bar{\F}_{\ell})
\]
都来自某个模形式。
\end{theorem}

\begin{remark}
这里“奇”是指复共轭元在表示下的行列式为 $-1$。该定理由 Khare--Wintenberger 等人完成证明，是现代数论的里程碑。
\end{remark}

\subsection*{Fermat 大定理的图景}

\begin{theorem}[模性定理，陈述]
定义在 $\Q$ 上的每条椭圆曲线都是模的。等价地，每条椭圆曲线都对应一个权 $2$ 的模形式，并且二者诱导的 Galois 表示相匹配。
\end{theorem}

\begin{proof}[与 Fermat 大定理的关系]
若 Fermat 方程
\[
a^p+b^p=c^p,
\qquad
p\ge3
\]
有非平凡整数解，则可构造 Frey 曲线
\[
E_{a,b,p}:y^2=x(x-a^p)(x+b^p)。
\]
Ribet 定理表明：由这样的解得到的 Frey 曲线会给出一个“不可能模”的模 $p$ Galois 表示；而 Wiles 及其后继者证明半稳定椭圆曲线都是模的，所以该表示又必须来自模形式。矛盾由此产生，因此 Fermat 方程没有非平凡整数解。
\end{proof}

\begin{remark}
从本章视角看，Fermat 大定理的证明链条是
\[
\text{整数方程}\to\text{椭圆曲线}\to\text{Galois 表示}\to\text{模形式}。
\]
Galois 理论在其中充当了连接算术、几何与分析的共同语言。
\end{remark}

\section*{习题}

本章共有 60 道习题，分为基础、进阶、挑战三级。

\subsection*{基础题}

\begin{enumerate}[label=\textbf{习题12.\arabic*.}, leftmargin=*, itemsep=0.75em]
\item \basic 设 $d\in\Z$ 为无平方因子整数且 $d\neq1$。证明 $\Gal(\Q(\sqrt d)/\Q)\cong\Z/2\Z$，并写出非平凡自同构。
\item \basic 计算 $\Gal(\Q(i)/\Q)$，并说明它与复共轭的关系。
\item \basic 设 $d_1,d_2$ 为互异无平方因子整数。求 $\Gal(\Q(\sqrt{d_1},\sqrt{d_2})/\Q)$。
\item \basic 求 $\Q(\sqrt2,\sqrt5)$ 的全部中间域。
\item \basic 对 $L=\Q(\sqrt2,\sqrt3)$ 画出中间域格与子群格。
\item \basic 求 $\alpha=\sqrt2+\sqrt3$ 在 $\Q$ 上的最小多项式，并确定其分裂域。
\item \basic 证明 $x^3-2$ 在 $\Q$ 上不可约，并求其分裂域与 Galois 群。
\item \basic 把上一题中的 $2$ 换成 $3$，重新计算分裂域与 Galois 群。
\item \basic 求 $x^3-1$ 的分裂域与 Galois 群。
\item \basic 计算 $\Gal(\F_9/\F_3)$，并写出其生成元。
\item \basic 计算 $\Gal(\F_{16}/\F_2)$，并说明它为何是循环群。
\item \basic 设 $n\mid m$。求 $\Gal(\F_{p^m}/\F_{p^n})$。
\item \basic 求 $\Gal(\Q(\zeta_3)/\Q)$。
\item \basic 求 $\Gal(\Q(\zeta_4)/\Q)$ 与 $\Gal(\Q(\zeta_6)/\Q)$。
\item \basic 计算 $\Gal(\Q(\zeta_5)/\Q)$，并找出一个生成元。
\item \basic 计算 $\Gal(\Q(\zeta_8)/\Q)$，并证明它不是循环群。
\item \basic 求 $\Q(\zeta_8)$ 的全部二次子域。
\item \basic 求 $\Phi_8(x)$, $\Phi_{10}(x)$, $\Phi_{12}(x)$。
\item \basic 对 $L=\Q(\sqrt2,\sqrt3)$，求 $\langle \sigma \rangle$、$\langle \tau \rangle$、$\langle \sigma\tau \rangle$ 的固定域。
\item \basic 设 $L=\F_{p^2}$。证明 $L^{\langle \Frob_p \rangle}=\F_p$。
\item \basic 证明任意 $\sigma\in\Gal(\Q(\zeta_n)/\Q)$ 都把原始 $n$ 次单位根送到另一个原始 $n$ 次单位根。
\item \basic 对 $L=\Q(\zeta_5)$，描述 $\sigma_2, \sigma_3, \sigma_4$ 的作用。
\item \basic 求 $x^4-2$ 的分裂域，并说明其 Galois 群的阶。
\item \basic 对 $L=\Q(i,\sqrt2)$，求 $\Gal(L/\Q)$ 并判断它是否循环。
\item \basic 证明 $\Q(\zeta_n)=\Q(\zeta_n^{-1})$，并指出复共轭对应于 $(\Z/n\Z)^\times$ 中哪个元素。
\end{enumerate}

\subsection*{进阶题}

\begin{enumerate}[resume, label=\textbf{习题12.\arabic*.}, leftmargin=*, itemsep=0.75em]
\item \medium 设 $L/K$ 为有限 Galois 扩张，$E$ 为中间域。证明 $[E:K]=[\Gal(L/K):\Gal(L/E)]$。
\item \medium 设 $L/K$ 为有限 Galois 扩张。证明 $E_1\subseteq E_2$ 当且仅当 $\Gal(L/E_2)\subseteq\Gal(L/E_1)$。
\item \medium 若 $\Gal(L/K)\cong C_p$，其中 $p$ 为素数，证明 $L/K$ 没有非平凡真中间域。
\item \medium 若 $\Gal(L/K)\cong C_{p^2}$，证明 $L/K$ 恰有一个次数为 $p$ 的中间域。
\item \medium 证明：中间域 $E/K$ 为 Galois 扩张当且仅当 $\Gal(L/E)\lhd\Gal(L/K)$。
\item \medium 求 $\Q(\zeta_{12})/\Q$ 的全部中间域，并写出对应子群。
\item \medium 利用分圆分解公式证明 $\sum_{d\mid n}\varphi(d)=n$。
\item \medium 证明对奇素数 $p$，有 $\Phi_{2p}(x)=\Phi_p(-x)$。
\item \medium 证明 $\Phi_{p^r}(x+1)$ 满足 Eisenstein 判别法，从而给出另一个不可约性证明。
\item \medium 证明 $x^{p^n}-x$ 等于所有次数整除 $n$ 的首一不可约多项式之积。
\item \medium 设 $n\mid m$。证明 $\Gal(\F_{p^m}/\F_{p^n})$ 由 $x\mapsto x^{p^n}$ 生成。
\item \medium 对 $p=7,11$，分别找出 $\Gal(\Q(\zeta_p)/\Q)$ 的一个生成元。
\item \medium 证明 $\Q(\zeta_p+\zeta_p^{-1})$ 是 $\Q(\zeta_p)$ 的最大实子域，且次数为 $(p-1)/2$。
\item \medium 利用 Galois 对应，分类 $\Q(\zeta_7)$ 的全部中间域。
\item \medium 设 $f(x)=x^3+ax+b$ 在 $\Q$ 上不可约，判别式为 $\Delta=-4a^3-27b^2$。证明其 Galois 群同构于 $A_3$ 当且仅当 $\Delta$ 是平方数。
\item \medium 把上一题应用到 $x^3-3x-1$ 与 $x^3-2$。
\item \medium 证明一切有限 $p$-群都是可解群。
\item \medium 证明二面体群 $D_{2n}$ 对任意 $n$ 都可解。
\item \medium 证明 $A_4$ 可解但不是阿贝尔群。
\item \medium 写出 $S_4$ 的一个阿贝尔商列，并据此解释一般四次方程为何可能由根式求解。
\item \medium 对 $L=\Q(\sqrt[3]{2},\omega)$，求所有中间域，并指出哪些是 $\Q$ 上的 Galois 扩张。
\item \medium 证明：若 $K$ 含有原始 $m$ 次单位根，则 $x^m-a$ 的分裂域 Galois 群是阿贝尔群。
\item \medium 证明 $\Gal(\Q(\zeta_9)/\Q)\cong C_6$，并分类其全部中间域。
\item \medium 计算 $\Gal(\Q(\zeta_{15})/\Q)$，并将其分解为两个循环群的直积。
\item \medium 解释正 $17$ 边形可作与 $\Gal(\Q(\zeta_{17})/\Q)$ 结构之间的关系。
\end{enumerate}

\subsection*{挑战题}

\begin{enumerate}[resume, label=\textbf{习题12.\arabic*.}, leftmargin=*, itemsep=0.75em]
\item \hard 证明 $x^5-2$ 在 $\Q$ 上的分裂域 Galois 群同构于 $C_5\rtimes C_4$，其阶为 $20$。
\item \hard 对上一题的分裂域，尽可能详细写出中间域与子群对应，至少给出所有指数为 $2,4,5$ 的子群所对应的固定域。
\item \hard 叙述逆 Galois 问题，并说明为什么分圆域为有限阿贝尔群的实现提供了重要线索。
\item \hard 设 $\rho:G_{\Q}\to\GL_2(E)$ 在 $p\neq\ell$ 处不分歧。证明若 $\Frob_p$ 与 $\Frob_p'$ 共轭，则它们在 $\rho$ 下的迹与行列式相同。
\item \hard 对有限 Galois 扩张与其一个复表示，写出 Artin $L$-函数的 Euler 乘积定义，并验证平凡表示时退化为 Dedekind $\zeta$-函数的局部因子。
\item \hard 设 $E/\Q$ 是椭圆曲线，$a_p=p+1-\#E(\F_p)$。说明为什么 $X^2-a_pX+p$ 可以看作某个二维 Galois 表示在 $\Frob_p$ 上的特征多项式。
\item \hard 详细说明 Wiles 证明 Fermat 大定理中的三步：Frey 曲线、Ribet 定理、模性定理，并指出每一步涉及的 Galois 表示。
\item \hard 设 $f$ 是 $\Q$ 上不可约五次多项式。证明若其 Galois 群作为 $S_5$ 的传递子群是可解的，则它嵌入 $C_5\rtimes C_4$。
\item \hard 设 $L/K$ 是 Galois 扩张且 $\Gal(L/K)\cong S_5$。证明不存在一个从 $1$ 到 $S_5$ 的正规列使所有商群都循环，并解释其与一般五次方程不可解性的关系。
\item \hard 设 $\bar\rho:G_{\Q}\to\GL_2(\bar{\F}_{\ell})$ 是连续不可约表示。解释什么叫奇表示，并说明来自椭圆曲线的模 $\ell$ 表示为什么是奇的。
\end{enumerate}

\section*{习题解答}

\begin{enumerate}[label=\textbf{\arabic*.}, leftmargin=*, itemsep=1em]
\item \textbf{解答.} 扩张 $\Q(\sqrt d)/\Q$ 的次数为 $2$，故任一 $\Q$-自同构由 $\sqrt d$ 的像唯一决定。又 $(\sigma(\sqrt d))^2=\sigma(d)=d$，所以 $\sigma(\sqrt d)=\pm\sqrt d$。恒等映射与
\[
\tau(a+b\sqrt d)=a-b\sqrt d\qquad(a,b\in\Q)
\]
都是 $\Q$-自同构，且后者非平凡，因此 $\Gal(\Q(\sqrt d)/\Q)=\{1,\tau\}\cong \Z/2\Z$。

\item \textbf{解答.} $\Q(i)$ 是 $x^2+1$ 的分裂域，$[\Q(i):\Q]=2$，故 Galois 群有两个元素。非平凡自同构必须把 $i$ 送到另一个根 $-i$，即
\[
\sigma(a+bi)=a-bi.
\]
这正是复共轭在 $\Q(i)$ 上的限制，所以 $\Gal(\Q(i)/\Q)\cong \Z/2\Z$，其非平凡元素就是复共轭。

\item \textbf{解答.} 设 $L=\Q(\sqrt{d_1},\sqrt{d_2})$。因 $d_1,d_2$ 为互异无平方因子整数，$\Q(\sqrt{d_1})\neq \Q(\sqrt{d_2})$，故 $[L:\Q]=4$。任一 $\Q$-自同构分别把 $\sqrt{d_1},\sqrt{d_2}$ 送到各自的相反数或本身，于是至多有四个自同构；四种符号选择都确实给出自同构。故
\[
\Gal(L/\Q)\cong (\Z/2\Z)^2,
\]
由“改变 $\sqrt{d_1}$ 符号”和“改变 $\sqrt{d_2}$ 符号”两个对合生成。

\item \textbf{解答.} 令 $L=\Q(\sqrt2,\sqrt5)$。它是双二次扩张，Galois 群为 Klein 四元群，故子群与中间域对应给出恰好三个非平凡真中间域。它们分别是
\[
\Q(\sqrt2),\qquad \Q(\sqrt5),\qquad \Q(\sqrt{10}).
\]
连同 $\Q$ 与 $L$ 本身，全部中间域就是
\[
\Q\subset \Q(\sqrt2),\ \Q(\sqrt5),\ \Q(\sqrt{10})\subset \Q(\sqrt2,\sqrt5).
\]

\item \textbf{解答.} 记 $L=\Q(\sqrt2,\sqrt3)$，令
\[
\sigma(\sqrt2)=-\sqrt2,\ \sigma(\sqrt3)=\sqrt3,\qquad
\tau(\sqrt2)=\sqrt2,\ \tau(\sqrt3)=-\sqrt3.
\]
则 $\Gal(L/\Q)=\{1,\sigma,\tau,\sigma\tau\}\cong V_4$。对应的中间域为
\[
L^{\langle\sigma\rangle}=\Q(\sqrt3),\quad L^{\langle\tau\rangle}=\Q(\sqrt2),\quad L^{\langle\sigma\tau\rangle}=\Q(\sqrt6).
\]
因此中间域格是顶部 $L$，中间三点 $\Q(\sqrt2),\Q(\sqrt3),\Q(\sqrt6)$，底部 $\Q$；子群格是顶部 $V_4$，中间三条阶 $2$ 子群 $\langle\sigma\rangle,\langle\tau\rangle,\langle\sigma\tau\rangle$，底部平凡群。两者按包含反向对应。

\item \textbf{解答.} 设 $\alpha=\sqrt2+\sqrt3$。则
\[
\alpha^2=5+2\sqrt6,
\]
故 $(\alpha^2-5)^2=24$，即
\[
\alpha^4-10\alpha^2+1=0.
\]
所以 $\alpha$ 满足 $f(x)=x^4-10x^2+1$。该多项式无有理根；若在 $\Q[x]$ 中分解为两个二次因子，则只能是
\[(x^2+ax-1)(x^2-ax-1)\quad\text{或}\quad (x^2+ax+1)(x^2-ax+1),\]
分别要求 $a^2=8$ 或 $a^2=12$，均不可能于 $\Q$ 中成立，故 $f$ 在 $\Q$ 上不可约。于是 $m_{\alpha,\Q}(x)=x^4-10x^2+1$。其根为
\[
\pm\sqrt2\pm\sqrt3,
\]
所以分裂域为 $\Q(\sqrt2,\sqrt3)$。

\item \textbf{解答.} 对 $f(x)=x^3-2$，由 Eisenstein 判别法（取素数 $2$）知它在 $\Q$ 上不可约，故 $[\Q(\sqrt[3]{2}):\Q]=3$。其三个根为
\[
\alpha=\sqrt[3]{2},\qquad \omega\alpha,\qquad \omega^2\alpha,
\]
其中 $\omega=(-1+\sqrt{-3})/2$。故分裂域为
\[
L=\Q(\alpha,\omega).
\]
因 $\Q(\alpha)\subset \R$ 而 $\Q(\omega)\not\subset \R$，故 $\Q(\alpha)\cap\Q(\omega)=\Q$，于是
\[
[L:\Q]=3\cdot2=6.
\]
Galois 群可由
\[
\sigma(\alpha)=\omega\alpha,\ \sigma(\omega)=\omega;
\qquad
\tau(\alpha)=\alpha,\ \tau(\omega)=\omega^2
\]
生成，满足 $\sigma^3=\tau^2=1$ 与 $\tau\sigma\tau=\sigma^{-1}$，故
\[
\Gal(L/\Q)\cong S_3.
\]

\item \textbf{解答.} 对 $f(x)=x^3-3$，同样由 Eisenstein 判别法（取素数 $3$）知其在 $\Q$ 上不可约。记 $\alpha=\sqrt[3]{3}$，则三个根为 $\alpha,\omega\alpha,\omega^2\alpha$，分裂域是
\[
L=\Q(\alpha,\omega).
\]
与上一题同理，$[L:\Q]=6$。定义
\[
\sigma(\alpha)=\omega\alpha,\ \sigma(\omega)=\omega,
\qquad
\tau(\alpha)=\alpha,\ \tau(\omega)=\omega^2.
\]
则它们生成六阶群并满足相同关系，因此
\[
\Gal(L/\Q)\cong S_3.
\]

\item \textbf{解答.} $x^3-1=(x-1)(x^2+x+1)$，其根为 $1,\omega,\omega^2$。故分裂域为 $\Q(\omega)=\Q(\sqrt{-3})$。因为 $x^2+x+1$ 在 $\Q$ 上不可约，所以
\[
[\Q(\omega):\Q]=2,
\]
故 Galois 群为二阶群，非平凡元素是复共轭 $\omega\mapsto\omega^2$。因此
\[
\Gal(x^3-1/\Q)=\Gal(\Q(\omega)/\Q)\cong \Z/2\Z.
\]

\item \textbf{解答.} $\F_9/\F_3$ 的次数为 $2$，有限域 Galois 群由 Frobenius 自同构生成，因此
\[
\Gal(\F_9/\F_3)=\langle \Frob_3\rangle,
\qquad \Frob_3(x)=x^3.
\]
又 $\Frob_3^2(x)=x^9=x$ 对一切 $x\in\F_9$ 成立，而 $\Frob_3\neq 1$，故其阶为 $2$。所以
\[
\Gal(\F_9/\F_3)\cong \Z/2\Z.
\]

\item \textbf{解答.} $\F_{16}/\F_2$ 的次数为 $4$。Frobenius 自同构
\[
\Frob_2:x\mapsto x^2
\]
固定 $\F_2$，且满足 $\Frob_2^4(x)=x^{16}=x$ 对所有 $x\in\F_{16}$ 成立。若 $\Frob_2^m=1$ 且 $m<4$，则一切元素都满足 $x^{2^m}=x$，从而 $X^{2^m}-X$ 有 $16$ 个根，矛盾。故 $\Frob_2$ 的阶为 $4$。于是
\[
\Gal(\F_{16}/\F_2)=\langle\Frob_2\rangle\cong \Z/4\Z,
\]
因而它是循环群。

\item \textbf{解答.} 若 $n\mid m$，则 $\F_{p^m}/\F_{p^n}$ 的扩张次数为 $m/n$。Galois 群由相对 Frobenius
\[
\varphi:x\mapsto x^{p^n}
\]
生成。的确，$\varphi$ 固定 $\F_{p^n}$，且
\[
\varphi^{m/n}(x)=x^{p^m}=x
\]
对所有 $x\in\F_{p^m}$ 成立；若 $0<r<m/n$ 就有 $p^{nr}<p^m$，故 $\varphi^r\neq 1$。因此
\[
\Gal(\F_{p^m}/\F_{p^n})\cong \Z/(m/n)\Z.
\]

\item \textbf{解答.} $\zeta_3=\omega$ 满足最小多项式 $x^2+x+1$，故 $[\Q(\zeta_3):\Q]=2$。唯一非平凡自同构由
\[
\sigma(\zeta_3)=\zeta_3^2
\]
给出，因此
\[
\Gal(\Q(\zeta_3)/\Q)\cong (\Z/3\Z)^\times\cong \Z/2\Z.
\]

\item \textbf{解答.} 由 $\zeta_4=i$，得 $\Q(\zeta_4)=\Q(i)$，故
\[
\Gal(\Q(\zeta_4)/\Q)\cong (\Z/4\Z)^\times=\{1,3\}\cong \Z/2\Z.
\]
又 $\zeta_6=e^{\pi i/3}$，而 $\zeta_6^2=\zeta_3$ 且 $\zeta_6=(1+\zeta_3)/2$，故 $\Q(\zeta_6)=\Q(\zeta_3)=\Q(\sqrt{-3})$，于是
\[
\Gal(\Q(\zeta_6)/\Q)\cong (\Z/6\Z)^\times=\{1,5\}\cong \Z/2\Z.
\]

\item \textbf{解答.} $\Phi_5(x)=x^4+x^3+x^2+x+1$，故
\[
[\Q(\zeta_5):\Q]=\varphi(5)=4.
\]
因此
\[
\Gal(\Q(\zeta_5)/\Q)\cong (\Z/5\Z)^\times.
\]
群 $(\Z/5\Z)^\times=\{1,2,3,4\}$ 是四阶循环群，例如 $2$ 的幂给出 $2,4,3,1$，故 $2$ 是生成元。于是
\[
\Gal(\Q(\zeta_5)/\Q)=\langle \sigma_2\rangle\cong C_4,
\qquad \sigma_2(\zeta_5)=\zeta_5^2.
\]

\item \textbf{解答.} 有
\[
\Gal(\Q(\zeta_8)/\Q)\cong (\Z/8\Z)^\times=\{1,3,5,7\}.
\]
其中 $3^2\equiv5^2\equiv7^2\equiv1\pmod8$，故每个非单位元素都阶为 $2$。因此
\[
(\Z/8\Z)^\times\cong C_2\times C_2,
\]
从而
\[
\Gal(\Q(\zeta_8)/\Q)\cong C_2\times C_2.
\]
它不是循环群，因为一个四阶循环群必须含有阶 $4$ 元素，而这里所有非平凡元素都只有阶 $2$。

\item \textbf{解答.} 由 $\zeta_8=(1+i)/\sqrt2$ 可知
\[
\Q(\zeta_8)=\Q(i,\sqrt2).
\]
其 Galois 群是 $V_4$，故有三个二次子域，分别对应三个阶 $2$ 子群。显然它们是
\[
\Q(i),\qquad \Q(\sqrt2),\qquad \Q(i\sqrt2)=\Q(\sqrt{-2}).
\]
这就是全部二次子域。

\item \textbf{解答.} 由分圆多项式定义或分解公式可得
\[
\Phi_8(x)=\frac{x^8-1}{x^4-1}=x^4+1,
\]
\[
\Phi_{10}(x)=\frac{x^{10}-1}{(x^5-1)(x^2-1)/(x-1)}=x^4-x^3+x^2-x+1,
\]
\[
\Phi_{12}(x)=\frac{x^{12}-1}{(x^6-1)(x^4-1)(x^3-1)(x^2-1)/(x-1)^2}=x^4-x^2+1.
\]
故答案为
\[
\Phi_8(x)=x^4+1,\qquad \Phi_{10}(x)=x^4-x^3+x^2-x+1,\qquad \Phi_{12}(x)=x^4-x^2+1.
\]

\item \textbf{解答.} 仍取
\[
\sigma(\sqrt2)=-\sqrt2,\ \sigma(\sqrt3)=\sqrt3,
\qquad
\tau(\sqrt2)=\sqrt2,\ \tau(\sqrt3)=-\sqrt3.
\]
则
\[
L^{\langle\sigma\rangle}=\Q(\sqrt3),
\qquad
L^{\langle\tau\rangle}=\Q(\sqrt2),
\qquad
L^{\langle\sigma\tau\rangle}=\Q(\sqrt6).
\]
理由是：$\sigma$ 只改变 $\sqrt2$ 的符号，故恰固定 $\Q(\sqrt3)$；$\tau$ 同理固定 $\Q(\sqrt2)$；而 $\sigma\tau$ 同时改变 $\sqrt2,\sqrt3$ 的符号，却固定它们的乘积 $\sqrt6$，故其固定域是 $\Q(\sqrt6)$。

\item \textbf{解答.} 设 $L=\F_{p^2}$，则 $\langle\Frob_p\rangle$ 的阶为 $2$。若 $x\in L^{\langle\Frob_p\rangle}$，则 $x^p=x$，故 $x$ 是多项式 $X^p-X$ 的根。这个多项式在任意特征 $p$ 域中恰有 $p$ 个根，即素域 $\F_p$ 的元素，所以 $x\in\F_p$。反过来，对任意 $a\in\F_p$ 都有 $a^p=a$，故 $\F_p\subseteq L^{\langle\Frob_p\rangle}$。因此
\[
L^{\langle\Frob_p\rangle}=\F_p.
\]

\item \textbf{解答.} 设 $\sigma\in\Gal(\Q(\zeta_n)/\Q)$，其中 $\zeta_n$ 是原始 $n$ 次单位根。因为 $\sigma$ 为域自同构，故保持乘法阶：若 $\zeta_n^m\neq1$，则 $\sigma(\zeta_n)^m=\sigma(\zeta_n^m)\neq1$；而 $\sigma(\zeta_n)^n=\sigma(1)=1$。所以 $\sigma(\zeta_n)$ 的阶仍是 $n$，即它仍是原始 $n$ 次单位根。等价地，必存在 $a\in(\Z/n\Z)^\times$ 使
\[
\sigma(\zeta_n)=\zeta_n^a.
\]

\item \textbf{解答.} 在 $L=\Q(\zeta_5)$ 中，定义 $\sigma_a(\zeta_5)=\zeta_5^a$。于是
\[
\sigma_2(\zeta_5)=\zeta_5^2,
\qquad
\sigma_3(\zeta_5)=\zeta_5^3=\zeta_5^{-2},
\qquad
\sigma_4(\zeta_5)=\zeta_5^4=\zeta_5^{-1}.
\]
其中 $\sigma_4$ 就是复共轭；又因为 $2\cdot3\equiv1\pmod5$，故 $\sigma_3=\sigma_2^{-1}$。在基 $1,\zeta_5,\zeta_5^2,\zeta_5^3$ 上，它们分别把 $\zeta_5^k$ 送到 $\zeta_5^{ak}$。

\item \textbf{解答.} 设 $\alpha=\sqrt[4]{2}$。多项式 $x^4-2$ 的根为
\[
\alpha,\ -\alpha,\ i\alpha,\ -i\alpha.
\]
故分裂域为
\[
L=\Q(\alpha,i).
\]
由 Eisenstein 判别法知 $x^4-2$ 在 $\Q$ 上不可约，所以 $[\Q(\alpha):\Q]=4$。而 $\Q(\alpha)\subset\R$，故 $i\notin\Q(\alpha)$，于是
\[
[L:\Q]=[L:\Q(\alpha)][\Q(\alpha):\Q]=2\cdot4=8.
\]
由于 $L/\Q$ 是分裂域，故为 Galois 扩张，从而
\[
|\Gal(L/\Q)|=[L:\Q]=8.
\]

\item \textbf{解答.} 设 $L=\Q(i,\sqrt2)$。这是两个不同二次域的复合域，故 $[L:\Q]=4$。任一 $\Q$-自同构独立地决定 $i\mapsto\pm i$、$\sqrt2\mapsto\pm\sqrt2$，所以 Galois 群恰有四个元素：
\[
1,\ c:i\mapsto-i,
\qquad s:\sqrt2\mapsto-\sqrt2,
\qquad cs.
\]
它们两两交换，且每个非平凡元素阶都为 $2$，故
\[
\Gal(L/\Q)\cong C_2\times C_2.
\]
因此它不是循环群。

\item \textbf{解答.} 因为 $\zeta_n^{-1}$ 已是原始 $n$ 次单位根，所以 $\Q(\zeta_n^{-1})\subseteq\Q(\zeta_n)$；反之 $\zeta_n=(\zeta_n^{-1})^{-1}\in\Q(\zeta_n^{-1})$，故
\[
\Q(\zeta_n)=\Q(\zeta_n^{-1}).
\]
复共轭把 $\zeta_n$ 送到 $\overline{\zeta_n}=\zeta_n^{-1}=\zeta_n^{n-1}$，因此在同构
\[
\Gal(\Q(\zeta_n)/\Q)\cong (\Z/n\Z)^\times
\]
下，复共轭对应于元素 $-1\pmod n$。

\item \textbf{解答.} 由有限 Galois 扩张的基本定理，设 $G=\Gal(L/K)$、$H=\Gal(L/E)$，则
\[
[L:E]=|H|,
\qquad [L:K]=|G|.
\]
利用次数乘法公式得
\[
[E:K]=\frac{[L:K]}{[L:E]}=\frac{|G|}{|H|}=[G:H]=[\Gal(L/K):\Gal(L/E)].
\]
证毕。

\item \textbf{解答.} 若 $E_1\subseteq E_2$，则任何固定 $E_2$ 的自同构当然也固定 $E_1$，故
\[
\Gal(L/E_2)\subseteq \Gal(L/E_1).
\]
反之若 $\Gal(L/E_2)\subseteq\Gal(L/E_1)$，取固定域得
\[
E_1=L^{\Gal(L/E_1)}\subseteq L^{\Gal(L/E_2)}=E_2.
\]
故两者等价。这正是 Galois 对应的反向包含性。

\item \textbf{解答.} 若 $\Gal(L/K)\cong C_p$，其中 $p$ 为素数，则它只有两个子群：平凡群和自身。由 Galois 对应，中间域与子群一一对应，所以只对应到底域 $K$ 与顶域 $L$。因此 $L/K$ 没有非平凡真中间域。

\item \textbf{解答.} 若 $\Gal(L/K)\cong C_{p^2}$，则循环群对每个整除其阶的正整数只有唯一子群。次数为 $p$ 的中间域 $E$ 对应满足
\[
[E:K]=p=[G:\Gal(L/E)]
\]
的子群，也就是 $G$ 中唯一的阶 $p$ 子群。故这样的中间域恰有一个。

\item \textbf{解答.} 设 $G=\Gal(L/K)$，$H=\Gal(L/E)$。若 $E/K$ 为 Galois 扩张，则对任意 $g\in G$ 与 $x\in E$，$g^{-1}(x)\in E$ 当且仅当 $x\in E$，故 $gHg^{-1}$ 固定 $E$。因此
\[
gHg^{-1}=\Gal(L/g(E))=\Gal(L/E)=H,
\]
所以 $H\lhd G$。反之若 $H\lhd G$，则对任意 $g\in G$ 有 $g(E)=L^{gHg^{-1}}=L^H=E$，故 $E$ 在 $G$ 作用下稳定，即 $E$ 是某个在 $K$ 上多项式根的共轭闭合集合生成的域，因此 $E/K$ 正规。又 $L/K$ 可分，故中间扩张 $E/K$ 也可分。于是 $E/K$ 为 Galois 扩张。

\item \textbf{解答.} 由
\[
(\Z/12\Z)^\times=\{1,5,7,11\}\cong C_2\times C_2
\]
知 $\Gal(\Q(\zeta_{12})/\Q)\cong V_4$。又
\[
\Q(\zeta_{12})=\Q(i,\sqrt3).
\]
因此中间域共有三个二次域：
\[
\Q(i),\qquad \Q(\sqrt3),\qquad \Q(\sqrt{-3})=\Q(\zeta_3).
\]
设 $\sigma_a(\zeta_{12})=\zeta_{12}^a$。因为 $i=\zeta_{12}^3$，$\sqrt3=\zeta_{12}+\zeta_{12}^{-1}$，可算得：
\[
\sigma_5: i\mapsto i,\ \sqrt3\mapsto-\sqrt3;
\qquad
\sigma_7: i\mapsto-i,\ \sqrt3\mapsto-\sqrt3;
\qquad
\sigma_{11}: i\mapsto-i,\ \sqrt3\mapsto\sqrt3.
\]
故对应关系为
\[
\Q(i)\leftrightarrow \langle\sigma_5\rangle,
\qquad
\Q(\sqrt{-3})\leftrightarrow \langle\sigma_7\rangle,
\qquad
\Q(\sqrt3)\leftrightarrow \langle\sigma_{11}\rangle.
\]
再加上 $\Q\leftrightarrow G$、$\Q(\zeta_{12})\leftrightarrow\{1\}$，即得全部中间域与子群。

\item \textbf{解答.} 分圆分解公式给出
\[
x^n-1=\prod_{d\mid n}\Phi_d(x).
\]
两边取次数，得到
\[
n=\deg(x^n-1)=\sum_{d\mid n}\deg\Phi_d(x)=\sum_{d\mid n}\varphi(d).
\]
故 $\sum_{d\mid n}\varphi(d)=n$。

\item \textbf{解答.} 设 $p$ 为奇素数。原始 $2p$ 次单位根恰是形如 $-\zeta$ 的数，其中 $\zeta$ 为原始 $p$ 次单位根；因为 $(-\zeta)^{2p}=1$，且其阶既不是 $1,2,p$，故为 $2p$。于是 $\Phi_{2p}(x)$ 的根集恰为 $\{-\zeta:\ \zeta\text{ 为原始 }p\text{ 次单位根}\}$。而 $\Phi_p(-x)$ 的根正是这些数，所以两者同为首一多项式且根集相同，故
\[
\Phi_{2p}(x)=\Phi_p(-x).
\]

\item \textbf{解答.} 设
\[
\Phi_{p^r}(x)=\frac{x^{p^r}-1}{x^{p^{r-1}}-1}=1+x^{p^{r-1}}+x^{2p^{r-1}}+\cdots+x^{(p-1)p^{r-1}}.
\]
令 $y=(x+1)^{p^{r-1}}-1$，则
\[
\Phi_{p^r}(x+1)=\frac{(1+y)^p-1}{y}=p+\binom p2 y+\cdots+y^{p-1}.
\]
注意到 $y$ 的各项系数都被 $p$ 整除，且常数项为 $0$；因此上式除首项外其余各项的每个系数都被 $p$ 整除。又常数项恰为 $p$，不被 $p^2$ 整除；最高次项系数为 $1$。所以 $\Phi_{p^r}(x+1)$ 满足 Eisenstein 判别法（素数为 $p$），从而它在 $\Q[x]$ 中不可约，故 $\Phi_{p^r}(x)$ 也不可约。

\item \textbf{解答.} 多项式 $x^{p^n}-x$ 的根正是 $\F_{p^n}$ 的全部元素。对任意 $\alpha\in\F_{p^n}$，其在 $\F_p$ 上的最小多项式次数等于 $[\F_p(\alpha):\F_p]$，而 $\F_p(\alpha)$ 是 $\F_{p^n}$ 的子域，因此该次数必整除 $n$。于是 $x^{p^n}-x$ 的每个不可约因子次数都整除 $n$。反过来，若 $f\in\F_p[x]$ 是次数整除 $n$ 的首一不可约多项式，设 $\deg f=d\mid n$，则其根属于 $\F_{p^d}\subseteq\F_{p^n}$，故 $f\mid x^{p^n}-x$。又
\[
(x^{p^n}-x)'=-1\neq0,
\]
所以没有重根。故
\[
x^{p^n}-x=\prod_{d\mid n}\ \prod_{\substack{f\ \text{首一不可约}\\ \deg f=d}} f(x).
\]

\item \textbf{解答.} 设 $m=nr$。上一道基础题已知 $\Gal(\F_{p^m}/\F_{p^n})$ 的阶为 $r$。映射
\[
\varphi:x\mapsto x^{p^n}
\]
固定 $\F_{p^n}$，故属于该 Galois 群。并且
\[
\varphi^r(x)=x^{p^{nr}}=x^{p^m}=x
\]
对一切 $x\in\F_{p^m}$ 成立。若 $0<t<r$ 且 $\varphi^t=1$，则 $x^{p^{nt}}=x$ 对所有 $x\in\F_{p^m}$ 成立，这将迫使 $\F_{p^m}\subseteq\F_{p^{nt}}$，与 $nt<m$ 矛盾。故 $\varphi$ 的阶恰为 $r= m/n$，因此它生成整个 Galois 群。

\item \textbf{解答.} 对素数 $p$，有
\[
\Gal(\Q(\zeta_p)/\Q)\cong (\Z/p\Z)^\times,
\]
这是阶 $p-1$ 的循环群。故只要找 $(\Z/p\Z)^\times$ 的原根即可。模 $7$ 下，$3$ 的幂为
\[
3,\ 3^2\equiv2,\ 3^3\equiv6,\ 3^4\equiv4,\ 3^5\equiv5,\ 3^6\equiv1 \pmod7,
\]
故 $3$ 生成 $(\Z/7\Z)^\times$，于是 $\sigma_3$ 生成 $\Gal(\Q(\zeta_7)/\Q)$。模 $11$ 下，$2$ 的幂为
\[
2,4,8,5,10,9,7,3,6,1 \pmod{11},
\]
故 $2$ 生成 $(\Z/11\Z)^\times$，于是 $\sigma_2$ 生成 $\Gal(\Q(\zeta_{11})/\Q)$。

\item \textbf{解答.} 设 $K=\Q(\zeta_p)$，$c$ 为复共轭，则 $c(\zeta_p)=\zeta_p^{-1}$。因此
\[
t:=\zeta_p+\zeta_p^{-1}
\]
被 $c$ 固定，故 $\Q(t)\subseteq K^{\langle c\rangle}$。由于 $c$ 在 $K$ 上非平凡且阶为 $2$，固定域次数为
\[
[K^{\langle c\rangle}:\Q]=\frac{[K:\Q]}{2}=\frac{p-1}{2}.
\]
另一方面，$K^{\langle c\rangle}$ 的元素都被复共轭固定，所以全是实数；于是它是 $K$ 中最大的实子域。又 $t$ 已在此固定域中，且其生成的域次数也必须是 $(p-1)/2$，故
\[
\Q(\zeta_p+\zeta_p^{-1})=K^{\langle c\rangle}
\]
就是最大实子域，其次数为 $(p-1)/2$。

\item \textbf{解答.} $\Gal(\Q(\zeta_7)/\Q)\cong (\Z/7\Z)^\times\cong C_6$，故其子群按阶唯一：阶 $1,2,3,6$ 各一个。由 Galois 对应，中间域也恰有四个：
\[
\Q\subset E_2, E_3\subset \Q(\zeta_7)\subset \Q(\zeta_7).
\]
具体地，唯一阶 $2$ 子群由复共轭 $\sigma_{-1}$ 生成，其固定域是唯一三次域，即最大实子域
\[
E_3=\Q(\zeta_7+\zeta_7^{-1}).
\]
唯一阶 $3$ 子群可取 $\langle\sigma_2\rangle$（因 $2^3\equiv1\pmod7$），其固定域是唯一二次域。对素数 $7\equiv3\pmod4$，该二次域为
\[
E_2=\Q(\sqrt{-7}).
\]
所以全部中间域为
\[
\Q,\qquad \Q(\sqrt{-7}),\qquad \Q(\zeta_7+\zeta_7^{-1}),\qquad \Q(\zeta_7).
\]

\item \textbf{解答.} 设 $f(x)=x^3+ax+b$ 的根为 $r_1,r_2,r_3$，其分裂域为 $L$。因 $f$ 在 $\Q$ 上不可约，Galois 群作为三根上的置换群是 $S_3$ 的传递子群，故只能是 $A_3$ 或 $S_3$。令
\[
\delta=(r_1-r_2)(r_1-r_3)(r_2-r_3),
\qquad \delta^2=\Delta.
\]
偶置换保持 $\delta$，奇置换改变 $\delta$ 的符号。因此 $\delta\in\Q$ 当且仅当 Galois 群中没有奇置换，也就是当且仅当 $\Gal(L/\Q)\subseteq A_3$。由于群只能是 $A_3$ 或 $S_3$，遂得
\[
\Gal(L/\Q)\cong A_3 \iff \delta\in\Q \iff \Delta=\delta^2\text{ 是 }\Q\text{ 中平方数}.
\]

\item \textbf{解答.} 对 $f(x)=x^3-3x-1$，有 $a=-3,b=-1$，故判别式
\[
\Delta=-4(-3)^3-27(-1)^2=108-27=81=9^2.
\]
由上一题得其 Galois 群同构于 $A_3\cong C_3$。对 $f(x)=x^3-2$，可写成 $x^3+0x-2$，故
\[
\Delta=-4\cdot0^3-27(-2)^2=-108,
\]
这不是有理平方数，因此其 Galois 群不是 $A_3$，只能是 $S_3$。

\item \textbf{解答.} 对有限 $p$-群 $G$ 用归纳法。若 $|G|=p$，则 $G$ 循环，当然可解。设结论对较小阶 $p$-群成立。由有限 $p$-群中心非平凡，取 $1\neq Z\le Z(G)$。则 $Z$ 为阿贝尔群，故可解；商群 $G/Z$ 仍是较小阶的 $p$-群，按归纳假设也可解。由“若正规子群与商群都可解，则整体可解”知 $G$ 可解。故一切有限 $p$-群都可解。

\item \textbf{解答.} 二面体群 $D_{2n}$ 由旋转 $r$ 与反射 $s$ 生成，满足 $r^n=s^2=1$、$srs=r^{-1}$。其旋转子群 $\langle r\rangle\cong C_n$ 是正规子群，且
\[
D_{2n}/\langle r\rangle\cong C_2.
\]
于是有正规列
\[
1\triangleleft \langle r\rangle \triangleleft D_{2n}
\]
其商群分别为 $C_n$ 与 $C_2$，都阿贝尔，所以 $D_{2n}$ 可解。

\item \textbf{解答.} $A_4$ 不是阿贝尔群，因为例如
\[
(123)(124)=(13)(24),\qquad (124)(123)=(14)(23),
\]
两者不同。另一方面，Klein 四元群
\[
V_4=\{1,(12)(34),(13)(24),(14)(23)\}
\]
是 $A_4$ 的正规子群，且
\[
1\triangleleft V_4\triangleleft A_4
\]
满足 $V_4/1\cong V_4$、$A_4/V_4\cong C_3$，两者都阿贝尔。故 $A_4$ 可解但不是阿贝尔群。

\item \textbf{解答.} 一个标准正规列为
\[
1\triangleleft V_4\triangleleft A_4\triangleleft S_4.
\]
其中
\[
V_4\cong C_2\times C_2,
\qquad A_4/V_4\cong C_3,
\qquad S_4/A_4\cong C_2.
\]
各商群都阿贝尔，所以 $S_4$ 可解。Galois 理论告诉我们：一般四次方程的 Galois 群嵌入 $S_4$，因而是可解群；而“Galois 群可解”正是方程可由根式求解的群论机制。这就是四次方程能有根式公式的根本原因。

\item \textbf{解答.} 设 $L=\Q(\sqrt[3]{2},\omega)$，其 Galois 群同构于 $S_3$。$S_3$ 的子群共有：平凡群、全群、唯一阶 $3$ 子群 $A_3=\langle\sigma\rangle$，以及三个阶 $2$ 子群。由 Galois 对应，全部中间域为：
\[
\Q,
\qquad L^{A_3}=\Q(\omega),
\qquad L^{\langle\tau\rangle}=\Q(\sqrt[3]{2}),
\qquad L^{\langle\sigma\tau\rangle}=\Q(\omega\sqrt[3]{2}),
\qquad L^{\langle\sigma^2\tau\rangle}=\Q(\omega^2\sqrt[3]{2}),
\qquad L.
\]
其中次数分别为 $1,2,3,3,3,6$。正规子群只有 $1,A_3,S_3$，故对应的 $\Q$ 上 Galois 中间扩张只有
\[
\Q,\qquad \Q(\omega),\qquad L.
\]
三个三次子域都不是 $\Q$ 上的 Galois 扩张，因为它们不是正规扩张。

\item \textbf{解答.} 设 $K$ 含有原始 $m$ 次单位根 $\zeta_m$，令 $\alpha^m=a$，则 $x^m-a$ 的全部根为 $\alpha,\zeta_m\alpha,\dots,\zeta_m^{m-1}\alpha$，都在 $K(\alpha)$ 中。所以分裂域就是 $K(\alpha)$。任一 $\sigma\in\Gal(K(\alpha)/K)$ 满足
\[
\sigma(\alpha)^m=\sigma(a)=a,
\]
故 $\sigma(\alpha)=\zeta_m^j\alpha$ 对某个 $j$ 成立。于是 $\sigma$ 完全由 $j$ 决定，并且对任意两自同构 $\sigma_j,\sigma_k$ 有
\[
\sigma_j\sigma_k(\alpha)=\zeta_m^{j+k}\alpha=\sigma_k\sigma_j(\alpha).
\]
故 Galois 群交换，即为阿贝尔群。事实上它嵌入单位根群 $\mu_m\cong C_m$。

\item \textbf{解答.} 由
\[
\Gal(\Q(\zeta_9)/\Q)\cong (\Z/9\Z)^\times
\]
且 $(\Z/9\Z)^\times=\{1,2,4,5,7,8\}$ 为六阶循环群（例如 $2$ 的阶为 $6$），故
\[
\Gal(\Q(\zeta_9)/\Q)\cong C_6.
\]
循环群 $C_6$ 对每个约数只有唯一子群，因此中间域恰有四个：
\[
\Q,
\qquad \text{唯一二次域},
\qquad \text{唯一三次域},
\qquad \Q(\zeta_9).
\]
其中唯一二次域对应唯一阶 $3$ 子群，即 $\Q(\zeta_3)=\Q(\sqrt{-3})$；唯一三次域对应唯一阶 $2$ 子群，即最大实子域
\[
\Q(\zeta_9+\zeta_9^{-1}).
\]
故全部中间域为
\[
\Q,\qquad \Q(\sqrt{-3}),\qquad \Q(\zeta_9+\zeta_9^{-1}),\qquad \Q(\zeta_9).
\]

\item \textbf{解答.} 由中国剩余定理，
\[
(\Z/15\Z)^\times\cong (\Z/3\Z)^\times\times(\Z/5\Z)^\times.
\]
右边分别同构于 $C_2$ 与 $C_4$，故
\[
\Gal(\Q(\zeta_{15})/\Q)\cong (\Z/15\Z)^\times\cong C_2\times C_4.
\]
例如 $2\pmod{15}$ 的阶为 $4$，$11\pmod{15}$ 的阶为 $2$，于是
\[
\Gal(\Q(\zeta_{15})/\Q)=\langle\sigma_2\rangle\times\langle\sigma_{11}\rangle\cong C_4\times C_2.
\]
这就是所求的两个循环群直积分解。

\item \textbf{解答.} Gauss--Wantzel 定理断言：正 $n$ 边形可尺规作当且仅当 $\varphi(n)$ 是 $2$ 的幂。对 $n=17$，有
\[
\varphi(17)=16=2^4,
\]
且
\[
\Gal(\Q(\zeta_{17})/\Q)\cong (\Z/17\Z)^\times\cong C_{16}.
\]
因此该 Galois 群的每个商层都可以取成二次扩张，等价地，$\Q(\zeta_{17})$ 的最大实子域 $\Q(\zeta_{17}+\zeta_{17}^{-1})$ 可由一串二次扩张构成。尺规作图正对应于不断开平方，所以正 $17$ 边形可作。换言之，$\Gal(\Q(\zeta_{17})/\Q)$ 是 $2$-群这一结构，正是可作性的群论原因。

\item \textbf{解答.} 设 $\alpha=\sqrt[5]{2}$，$\zeta=\zeta_5$。由 Eisenstein 判别法，$x^5-2$ 在 $\Q$ 上不可约，故 $[\Q(\alpha):\Q]=5$。其分裂域为
\[
L=\Q(\alpha,\zeta),
\]
因为五个根是 $\alpha,\zeta\alpha,\zeta^2\alpha,\zeta^3\alpha,\zeta^4\alpha$。又 $[\Q(\zeta):\Q]=\varphi(5)=4$，且 $\Q(\alpha)\subset\R$、$\Q(\zeta)$ 没有非平凡实子扩张与次数 $5$ 相交，故
\[
\Q(\alpha)\cap\Q(\zeta)=\Q.
\]
于是
\[
[L:\Q]=5\cdot4=20.
\]
定义自同构
\[
\sigma(\alpha)=\zeta\alpha,\ \sigma(\zeta)=\zeta;
\qquad
\tau(\alpha)=\alpha,\ \tau(\zeta)=\zeta^2.
\]
则 $\sigma^5=1$，$\tau^4=1$，且
\[
\tau\sigma\tau^{-1}(\alpha)=\tau\sigma(\alpha)=\tau(\zeta\alpha)=\zeta^2\alpha=\sigma^2(\alpha),
\]
故 $\tau\sigma\tau^{-1}=\sigma^2$。因此 $\langle\sigma\rangle\cong C_5$ 为正规子群，$\langle\tau\rangle\cong C_4$ 通过自同构 $x\mapsto x^2$ 作用于它，从而
\[
\Gal(L/\Q)\cong C_5\rtimes C_4,
\]
其阶为 $20$。

\item \textbf{解答.} 仍记 $G=\langle\sigma,\tau\mid \sigma^5=\tau^4=1,\ \tau\sigma\tau^{-1}=\sigma^2\rangle$，$L=\Q(\alpha,\zeta_5)$，其中 $\alpha=\sqrt[5]{2}$。Galois 对应把子群与固定域反向对应。

先给出最重要的几类子群。
\begin{itemize}
\item 指数 $4$ 子群即阶 $5$ 子群：Sylow $5$ 子群唯一，等于 $\langle\sigma\rangle$。其固定域为
\[
L^{\langle\sigma\rangle}=\Q(\zeta_5),
\]
因为 $\sigma$ 固定 $\zeta_5$ 而循环置换五个根。
\item 指数 $5$ 子群即阶 $4$ 子群：共有 $5$ 个，分别是
\[
H_k=\langle \sigma^k\tau\rangle\qquad (k=0,1,2,3,4).
\]
令 $\beta_k=\zeta_5^{-k}\alpha$，则
\[
(\sigma^k\tau)(\beta_k)=\sigma^k(\tau(\zeta_5^{-k}\alpha))=\sigma^k(\zeta_5^{-2k}\alpha)=\zeta_5^{-k}\alpha=\beta_k.
\]
故
\[
L^{H_k}=\Q(\beta_k)=\Q(\zeta_5^{-k}\sqrt[5]{2}).
\]
这是五个互为共轭的五次子域；其中 $H_0=\langle\tau\rangle$ 的固定域是 $\Q(\sqrt[5]{2})$。
\item 指数 $2$ 子群即阶 $10$ 子群：唯一这样的子群是
\[
N=\langle\sigma,\tau^2\rangle.
\]
这里 $\tau^2$ 在 $\Q(\zeta_5)$ 上是复共轭，故
\[
L^N=\Q(\zeta_5)^{\langle\tau^2\rangle}=\Q(\zeta_5+\zeta_5^{-1})=\Q(\sqrt5).
\]
\end{itemize}

再补充其余子群。共有 $5$ 个阶 $2$ 子群
\[
J_k=\langle \sigma^k\tau^2\rangle\qquad (k=0,1,2,3,4),
\]
其固定域次数为 $10$。记 $t=\zeta_5+\zeta_5^{-1}$，取 $\gamma_k=\zeta_5^{3k}\alpha$，则 $J_k$ 固定 $t$ 与 $\gamma_k$，故可写作
\[
L^{J_k}=\Q(t,\gamma_k)=\Q(\zeta_5+\zeta_5^{-1},\ \zeta_5^{3k}\sqrt[5]{2}).
\]
连同 $\Q\leftrightarrow G$、$L\leftrightarrow\{1\}$，这就给出了全部中间域与子群对应。特别地，指数 $2,4,5$ 的情形分别对应于 $\Q(\sqrt5)$、$\Q(\zeta_5)$ 与五个五次子域 $\Q(\zeta_5^{-k}\sqrt[5]{2})$。

\item \textbf{解答.} 逆 Galois 问题问：给定任意有限群 $G$，是否存在一个 $\Q$ 的有限 Galois 扩张 $L/\Q$ 使
\[
\Gal(L/\Q)\cong G?
\]
这是数论中的核心公开问题之一。分圆域之所以提供重要线索，是因为
\[
\Gal(\Q(\zeta_n)/\Q)\cong (\Z/n\Z)^\times
\]
总是阿贝尔群，而且 Kronecker--Weber 定理进一步说明：每个有限阿贝尔扩张都包含在某个分圆域中。反过来，任意有限阿贝尔群都可通过适当选择 $n$ 并取分圆域中的合适固定子域来实现。因此分圆域完全解决了阿贝尔情形，也为“通过构造特殊多项式和特殊扩张来实现给定群”提供了最基本的范例。

\item \textbf{解答.} “在 $p\neq\ell$ 处不分歧”意味着 $\rho$ 在惯性子群 $I_p$ 上平凡，因此 Frobenius 元只在共轭意义下定义。若 $\Frob_p'$ 与 $\Frob_p$ 共轭，即存在 $g\in G_{\Q}$ 使
\[
\Frob_p'=g\Frob_p g^{-1},
\]
则由表示同态性得
\[
\rho(\Frob_p')=\rho(g)\rho(\Frob_p)\rho(g)^{-1}.
\]
矩阵的迹与行列式在共轭下不变，因此
\[
\tr\rho(\Frob_p')=\tr\rho(\Frob_p),
\qquad
\det\rho(\Frob_p')=\det\rho(\Frob_p).
\]
所以迹与行列式只依赖于 Frobenius 共轭类，而不依赖于代表元的选择。

\item \textbf{解答.} 设 $L/K$ 为有限 Galois 扩张，$G=\Gal(L/K)$，$\rho:G\to\GL(V)$ 为一个复表示。对 $K$ 的每个非零素理想 $\mathfrak p$，记 $I_{\mathfrak p}$ 为惯性群，$\Frob_{\mathfrak p}$ 为几何 Frobenius 共轭类，则 Artin $L$-函数定义为
\[
L(s,\rho)=\prod_{\mathfrak p}\det\Bigl(1-\rho(\Frob_{\mathfrak p})\,N\mathfrak p^{-s}\ \big|\ V^{I_{\mathfrak p}}\Bigr)^{-1}.
\]
若 $\rho=\mathbf 1$ 是平凡表示，则对每个 $\mathfrak p$ 都有 $V^{I_{\mathfrak p}}=V\cong\C$，且 $\rho(\Frob_{\mathfrak p})=1$，所以局部因子化为
\[
\det(1-N\mathfrak p^{-s})^{-1}=(1-N\mathfrak p^{-s})^{-1}.
\]
这恰是 Dedekind $\zeta$-函数 $\zeta_K(s)$ 的 Euler 因子，因此
\[
L(s,\mathbf1)=\zeta_K(s).
\]

\item \textbf{解答.} 对椭圆曲线 $E/\Q$ 与素数 $\ell\neq p$，其 $\ell$-进 Tate 模给出二维表示
\[
\rho_{E,\ell}:G_{\Q}\to\GL_2(\Z_\ell).
\]
当 $p$ 是 $E$ 的良素数时，$\rho_{E,\ell}$ 在 $p$ 处不分歧，且 Weil 猜想（在椭圆曲线情形由 Hasse 与 Grothendieck 理论实现）给出
\[
\det\bigl(1-X\rho_{E,\ell}(\Frob_p)\bigr)=1-a_pX+pX^2,
\]
其中
\[
a_p=p+1-\#E(\F_p).
\]
把变量改写成特征多项式变量 $T$，就得到
\[
\chi_{\rho_{E,\ell}(\Frob_p)}(T)=T^2-a_pT+p.
\]
因此 $X^2-a_pX+p$ 正是该二维 Galois 表示在 $\Frob_p$ 上的特征多项式。

\item \textbf{解答.} Wiles 证明 Fermat 大定理的逻辑链可以概括为三步。
\begin{enumerate}[label=\normalfont(\arabic*), leftmargin=2.5em]
\item 若存在非平凡解 $a^p+b^p=c^p$（$p>2$），则构造 Frey 曲线
\[
E_{a,b,p}: y^2=x(x-a^p)(x+b^p).
\]
这条曲线是半稳定的，并带来一个模 $p$ Galois 表示
\[
\bar\rho_{E,p}:G_{\Q}\to\GL_2(\F_p)
\]
来自其 $p$-挠点 $E[p]$。
\item Ribet 定理（即 epsilon 猜想）说明：若这样的 Frey 曲线存在，则其模 $p$ 表示虽然应当是模的，却会被降层到极小的 level $2$。更准确地说，若 $\bar\rho_{E,p}$ 来自某个模形式，则它实际上来自权 $2$、level $2$ 的模形式。但这样的模形式不存在，因此得到矛盾。这里核心对象仍是 $\bar\rho_{E,p}$。
\item Wiles 的模性定理（后由 Taylor--Wiles 完成推广）证明所有半稳定椭圆曲线都是模的。于是 Frey 曲线必然模，从而其 $\ell$-进表示
\[
\rho_{E,\ell}:G_{\Q}\to\GL_2(\Z_\ell)
\]
以及相应的模 $p$ 表示都来自模形式。与 Ribet 的降层结论结合，即推出 Frey 曲线不可能存在，因此 Fermat 方程无非平凡解。
\end{enumerate}
所以三步中涉及的表示分别是：Frey 曲线的模 $p$ 表示、Ribet 降层处理的同一模 $p$ 表示，以及 Wiles 比较的椭圆曲线 $\ell$-进表示与模形式附带的 Galois 表示。

\item \textbf{解答.} 设 $G\le S_5$ 是 $f$ 的 Galois 群在五个根上的置换像。因为 $f$ 不可约，$G$ 在 $5$ 个根上传递。又假设 $G$ 可解。对素数次数 $5$ 的传递可解置换群，存在一个正规 Sylow $5$ 子群 $P$，并且 $P$ 的作用是正则的；这是 Burnside 关于素数次数可解传递群的结论。于是 $P\cong C_5$，而点稳定子 $G_\alpha$ 与 $P$ 交平凡，且
\[
G=P\rtimes G_\alpha.
\]
因为 $G_\alpha$ 通过共轭作用忠实嵌入 $\Aut(P)$，而
\[
\Aut(C_5)\cong (\Z/5\Z)^\times\cong C_4,
\]
故 $G_\alpha\le C_4$。因此
\[
G\hookrightarrow C_5\rtimes C_4.
\]
这就证明了任何可解的五次不可约多项式的传递 Galois 群都嵌入该 Frobenius 群。

\item \textbf{解答.} 若存在正规列
\[
1=G_0\triangleleft G_1\triangleleft\cdots\triangleleft G_r=S_5
\]
使每个商群 $G_{i+1}/G_i$ 都循环，则这些商群都阿贝尔，从而 $S_5$ 为可解群。但 $S_5$ 不是可解群：因为它有正规子群 $A_5$，而 $A_5$ 是非阿贝尔单群，所以在任何正规列中一旦出现 $A_5$，其再往下的商不可能分解为循环群；尤其 $A_5$ 本身不是阿贝尔群。故这样的正规列不存在。

这与一般五次方程不可由根式求解的关系是：一个多项式若能由根式求解，则其分裂域 Galois 群必须可解。一般五次方程的“泛”Galois 群是 $S_5$，而 $S_5$ 不可解，所以不存在统一的根式公式来解一般五次方程。

\item \textbf{解答.} 对二维表示 $\bar\rho:G_{\Q}\to\GL_2(\bar\F_\ell)$，称它是\emph{奇}的，如果对任一复共轭元 $c\in G_{\Q}$ 有
\[
\det\bar\rho(c)=-1.
\]
当 $\ell\neq2$ 时，这等价于 $\bar\rho(c)$ 的特征值为 $1$ 与 $-1$，也等价于 $\bar\rho(c)$ 不为标量矩阵。若行列式为 $+1$，则称为偶表示。

对椭圆曲线 $E/\Q$，其模 $\ell$ 表示
\[
\bar\rho_{E,\ell}:G_{\Q}\to\GL_2(E[\ell])\cong\GL_2(\F_\ell)
\]
来自 $\ell$-挠点。复共轭在 $E(\C)$ 上诱导一个反全纯作用；在线性化后，它在 Tate 模（从而在 $E[\ell]$）上的特征值是 $1$ 与 $-1$，所以
\[
\det\bar\rho_{E,\ell}(c)=-1.
\]
因此来自椭圆曲线的模 $\ell$ 表示是奇表示。
\end{enumerate}

